% Options for packages loaded elsewhere
\PassOptionsToPackage{unicode}{hyperref}
\PassOptionsToPackage{hyphens}{url}
\PassOptionsToPackage{dvipsnames,svgnames,x11names}{xcolor}
\documentclass[
  11pt,
]{article}
\usepackage{xcolor}
\usepackage[margin=0.85in]{geometry}
\usepackage{amsmath,amssymb}
\setcounter{secnumdepth}{-\maxdimen} % remove section numbering
\usepackage{iftex}
\ifPDFTeX
  \usepackage[T1]{fontenc}
  \usepackage[utf8]{inputenc}
  \usepackage{textcomp} % provide euro and other symbols
\else % if luatex or xetex
  \usepackage{unicode-math} % this also loads fontspec
  \defaultfontfeatures{Scale=MatchLowercase}
  \defaultfontfeatures[\rmfamily]{Ligatures=TeX,Scale=1}
\fi
\usepackage{lmodern}
\ifPDFTeX\else
  % xetex/luatex font selection
\fi
% Use upquote if available, for straight quotes in verbatim environments
\IfFileExists{upquote.sty}{\usepackage{upquote}}{}
\IfFileExists{microtype.sty}{% use microtype if available
  \usepackage[]{microtype}
  \UseMicrotypeSet[protrusion]{basicmath} % disable protrusion for tt fonts
}{}
\makeatletter
\@ifundefined{KOMAClassName}{% if non-KOMA class
  \IfFileExists{parskip.sty}{%
    \usepackage{parskip}
  }{% else
    \setlength{\parindent}{0pt}
    \setlength{\parskip}{6pt plus 2pt minus 1pt}}
}{% if KOMA class
  \KOMAoptions{parskip=half}}
\makeatother
\setlength{\emergencystretch}{3em} % prevent overfull lines
\providecommand{\tightlist}{%
  \setlength{\itemsep}{0pt}\setlength{\parskip}{0pt}}
\usepackage{amsmath,amssymb}
\usepackage{float}
\usepackage{microtype}
\usepackage{fvextra}
\fvset{breaklines=true}
\usepackage{setspace}
\usepackage{lineno}
\usepackage{etoolbox}
\usepackage{needspace}
\usepackage{xurl}
\AtBeginEnvironment{longtable}{\singlespacing\small}
\AtBeginEnvironment{verbatim}{\singlespacing\small}
\setlength{\emergencystretch}{3em}
\floatplacement{figure}{H}
\DeclareUnicodeCharacter{2212}{\ensuremath{-}}
\DeclareUnicodeCharacter{2264}{\ensuremath{\leq}}
\DeclareUnicodeCharacter{2265}{\ensuremath{\geq}}
\DeclareUnicodeCharacter{2260}{\ensuremath{\ne}}
\DeclareUnicodeCharacter{2208}{\ensuremath{\in}}
\DeclareUnicodeCharacter{221E}{\ensuremath{\infty}}
\DeclareUnicodeCharacter{03B4}{\ensuremath{\delta}}
\DeclareUnicodeCharacter{03BB}{\ensuremath{\lambda}}
\DeclareUnicodeCharacter{03C4}{\ensuremath{\tau}}
\DeclareUnicodeCharacter{03BA}{\ensuremath{\kappa}}
\DeclareUnicodeCharacter{03B3}{\ensuremath{\gamma}}
\DeclareUnicodeCharacter{03B5}{\ensuremath{\varepsilon}}
\DeclareUnicodeCharacter{03C6}{\ensuremath{\phi}}
\DeclareUnicodeCharacter{2124}{\ensuremath{\mathbb Z}}
\DeclareUnicodeCharacter{1D63}{\ensuremath{{}_r}}
\DeclareUnicodeCharacter{2080}{\ensuremath{{}_0}}
\DeclareUnicodeCharacter{208A}{\ensuremath{{}_+}}
\DeclareUnicodeCharacter{208B}{\ensuremath{{}_-}}
\DeclareUnicodeCharacter{207B}{\ensuremath{{}^{-}}}
\DeclareUnicodeCharacter{00B9}{\ensuremath{{}^{1}}}
\DeclareUnicodeCharacter{00B2}{\ensuremath{{}^{2}}}
\DeclareUnicodeCharacter{2074}{\ensuremath{{}^{4}}}
\DeclareUnicodeCharacter{2075}{\ensuremath{{}^{5}}}
\DeclareUnicodeCharacter{2076}{\ensuremath{{}^{6}}}
\DeclareUnicodeCharacter{2192}{\ensuremath{\rightarrow}}
\DeclareUnicodeCharacter{0108}{\ensuremath{\widehat C}}
\DeclareUnicodeCharacter{00C2}{\ensuremath{\widehat A}}

\AtBeginDocument{\setstretch{1.12}}
\usepackage{bookmark}
\IfFileExists{xurl.sty}{\usepackage{xurl}}{} % add URL line breaks if available
\urlstyle{same}
% fallback for those not using the hyperref driver hyperxmp:
\makeatletter
\@ifundefined{xmpquote}{\newcommand{\xmpquote}[1]{#1}}{}
\makeatother
\hypersetup{
  colorlinks=true,
  linkcolor={Maroon},
  filecolor={Maroon},
  citecolor={Blue},
  urlcolor={blue},
  pdfcreator={LaTeX via pandoc}}

\author{}
\date{}

\begin{document}

\section{Additional file 2: Anchor
theory}\label{additional-file-2-anchor-theory}

\subsection{1. Object of study and
interpretation}\label{object-of-study-and-interpretation}

Let \(D\) and \(E\) be finite persistence diagrams with finite
coordinates, including the diagonal with unlimited multiplicity in their
bottleneck comparison. Write \(E+t\) for translation of both coordinates
by \(t\) and define

\[
f(t)=d_B(D,E+t).
\]

For the present superlevel contact construction, implemented as a
sublevel filtration of \(-L\), subtracting a map mean from \(L\) adds
that mean to its persistence coordinates. Consequently, with
\(\delta=\operatorname{mean}(L_{alt})-\operatorname{mean}(L_{ref})\),
the native and mean-centered distances are \(f(0)\) and \(f(\delta)\).
Common translation of both diagrams cancels. Every individual interval
lifetime is unchanged by this operation.

This identity does not say that a centered score is a function of the
native scalar distance and mean change alone. For \(D=\{(0,4)\}\),
\(E_+=\{(1,5)\}\) and \(E_-=\{(-1,3)\}\), both native distances equal 1.
Under the same \(\delta=1/4\) their distances become 5/4 and 3/4. The
direction is information about diagram location lost in the native
scalar summary. Appropriate interpretation: \textbf{anchor-location
sensitivity of a diagram comparison}. Inappropriate interpretation:
proof that mean intensity alone explains a biological effect, or that
centered diagrams lack structural information.

\subsection{2. Endpoint rigidity and an error-aware path
tube}\label{endpoint-rigidity-and-an-error-aware-path-tube}

Translation and the triangle inequality give \(|f(t)-f(s)|\leq |t-s|\).
This elementary inequality yields a stronger diagnostic than checking
selected intermediate shifts.

Let \(\delta\ne 0\), \(T=|\delta|\), \(h=\operatorname{sign}(\delta)\),
and \(g(u)=f(hu)\), \(0\leq u\leq T\). Put \(A=g(0)\), \(C=g(T)\). For
every \(u\),

\[
\max\{A-u,C-(T-u)\}\leq g(u)\leq\min\{A+u,C+(T-u)\}.\qquad (1)
\]

For either \(r\in \{-1,+1\}\), define the saturation deficit

\[
\kappa_r=T-r(C-A)\geq0.
\]

Then, throughout the entire segment,

\[
u-\kappa_r\leq r[g(u)-A]\leq u.\qquad (2)
\]

Proof. Apply the 1-Lipschitz inequality from the starting endpoint and
from the ending endpoint, obtaining (1); substitute
\(C-A=r(T-\kappa_r)\). Both signs give (2). No assumption about a fixed
optimal matching is needed.

If \(|C-A|=T\), taking \(r=\operatorname{sign}(C-A)\) gives
\(\kappa_r=0\) and therefore \(g(u)=A+ru\) everywhere on the segment.
This is an exact theorem for any shift magnitude, not just an
infinitesimal result. It does \textbf{not} imply that one matching
remains optimal, or that this affine formula extends outside the
segment.

Suppose only estimates \(\widehat A\), \(\widehat C\) are available,
with proven bounds \(|\widehat A-A|\leq e_0\) and
\(|\widehat C-C|\leq e_T\). Then

\[
\kappa_r\leq K_r:=T-r(\widehat C-\widehat A)+e_0+e_T.\qquad (3)
\]

For valid error bounds \(K_r\) is nonnegative. Combining (2) and (3),
the whole path lies within \(e_0+K_r\) of the line \(\widehat A+ru\). A
sharper signed statement is \(r[g(u)-A]\in [u-K_r,u]\). For nonsaturated
paths, the following envelopes need no affine approximation:

\[
\begin{aligned}L(u)&=\max\{0,\widehat A-e_0-u,\widehat C-e_T-(T-u)\},\\ U(u)&=\min\{\widehat A+e_0+u,\widehat C+e_T+(T-u)\}.\end{aligned}
\]

They satisfy \(L(u)\leq g(u)\leq U(u)\). At \(\delta=0\) there is no
path to audit.

There are two distinct computational claims. A residual at machine
precision is \textbf{observed numerical saturation}. A certified \(K\)
requires endpoint error bounds covering every relevant operation: PH
input conversion, coordinate recovery and translation, mean estimation
when an exact-real mean is the target, distance evaluation and its
returned rounding. The existing float32 filtration audit alone does not
verify all binary64 arithmetic. One can conservatively include such
errors, but must state the budget and its provenance. A solver's exact
algorithm mode is not, by itself, a formal interval-arithmetic proof for
a real-valued source map. API quantization remains part of the delivered
data and cannot be reversed by this audit.

If \(\delta\) is uncertain by \(e_\delta\), either define the audit
target as the explicitly stored shift, or add \(e_\delta\) to the
endpoint/path error when comparing to the intended shift; 1-Lipschitz
continuity justifies that addition.

\subsection{3. Why local slopes have only three possible
values}\label{why-local-slopes-have-only-three-possible-values}

For \(x=(b_x,d_x)\in D\) and \(y=(b_y,d_y)\in E\), set

\[
\begin{aligned}c_{xy}&=((b_x-b_y)+(d_x-d_y))/2,\\ q_{xy}&=|(b_x-b_y)-(d_x-d_y)|/2.\end{aligned}
\]

Their matching cost is exactly

\[
\varphi_{xy}(t)=\max\{|b_x-b_y-t|,|d_x-d_y-t|\}=|t-c_{xy}|+q_{xy}.
\]

Matching a point to the diagonal costs half its lifetime, independently
of \(t\). For any fixed partial matching \(M\), the maximum cost
therefore has the form \(\max\{a_M-t,b_M+t,c_M\}\), allowing absent
real-real terms and taking \(c_M\geq 0\). There are finitely many
partial matchings. Thus \(f\) is a finite lower envelope of these
functions: continuous, piecewise affine, 1-Lipschitz, with affine slopes
in \(\{-1,0,+1\}\). It need not be convex: for \(D=E=\{(0,2)\}\),
\(f(t)=\min\{|t|,1\}\).

The familiar parameterized-matching structure is a useful derivation for
this audit, not a claim of new mathematical machinery.

\subsubsection{A native-diagram-only local
certificate}\label{a-native-diagram-only-local-certificate}

The following provides a prospective slope from native diagrams, without
consulting a centered endpoint or selected-allele labels. If either
finite diagram is empty, \(f\) is constant under translation, with slope
zero for every shift. Otherwise, assume \(b=f(0)>0\). Form the standard
finite augmented bipartite graph: real \(D\) points and \(E\)-diagonal
copies on the left; real \(E\) points and \(D\)-diagonal copies on the
right. Include every real-real edge, each real point's own diagonal
edge, and all diagonal-copy to diagonal-copy edges of cost zero. A
perfect matching encodes a persistence matching.

Call an edge low, critical or high according as its native cost is
\(<b\), \(=b\) or \(>b\). For a direction \(h\in \{-1,+1\}\), a critical
real-real edge has one-sided slope \(a_e=-1\) if \(h c_{xy}>0\), and
\(+1\) otherwise. Critical diagonal edges have slope zero. Define

\[
\begin{aligned}\gamma&=\min\{|w_e(0)-b|:w_e(0)\ne b\},\\ R&=\min\{\gamma/2,\ |c_{xy}|\text{ for critical real-real edges with }h c_{xy}>0\}.\end{aligned}
\]

An empty second minimum is \(+\infty\). Zero-cost dummy edges ensure the
first minimum exists when \(b>0\). For \(0<u<R\), low edges stay below
each candidate level \(b+ru\), high edges stay above, and critical edges
retain their one-sided affine formulas.

Let \(G_r\) include all low edges and the critical edges with
\(a_e\leq r\). The smallest \(r\in \{-1,0,+1\}\) for which \(G_r\) has a
perfect matching satisfies

\[
f(hu)=b+ru\quad\text{for all }0\leq u<R.\qquad (4)
\]

Proof. The stated inequalities fix which edges can be used at the
threshold \(b+ru\). A native-optimal matching must contain at least one
critical edge, since otherwise its native maximum would be \(<b\).
Therefore the minimum, over optimal matchings, of their largest critical
slope is the exact active slope; nonoptimal matchings cannot enter below
it before the gap closes. This is equivalently the three
threshold-matching tests above.

This radius is conservative and may be much smaller than an empirical
mean shift. Predicting the native one-sided slope does not justify
extending it beyond \(R\). Cases \(b=0\) should use a separate
degeneracy treatment, rather than apply this statement unchanged. The
certificate can be implemented with exact dyadic arithmetic for stored
binary coordinates, but dense graph construction is optional and
potentially expensive.

If all coordinates in \(D\) and \(E\) lie on one affine lattice
\(a+g\mathbb Z\), then real-real costs are multiples of \(g\) and
diagonal costs are multiples of \(g/2\). Hence \(\gamma\geq g/2\) and
every positive \(|c_{xy}|\geq g/2\), so \(R\geq g/4\). This is a
conservative lattice-derived local radius, not an assertion that an
observed shift lies within it. Source quantization alone is insufficient
to choose \(g\) after diagram-coordinate recovery and translation; the
actual stored coordinates must be checked. Small shifts on a nonflat
affine cell are necessarily saturated. The empirical frequency of
saturation should therefore be compared with shift magnitudes,
candidate-level gaps and flat/turning synthetic controls, rather than
treated as a novel phenomenon.

\subsection{4. Witnesses that can avoid another complete matching
computation}\label{witnesses-that-can-avoid-another-complete-matching-computation}

A diagonal-augmented Hausdorff lower bound is

\[
\begin{aligned}H(t)=\max\bigl\{&\max_x\min[\operatorname{dist}_{\infty}(x,E+t),\operatorname{life}(x)/2],\\ &\max_y\min[\operatorname{dist}_{\infty}(y+t,D),\operatorname{life}(y)/2]\bigr\}\leq f(t).\end{aligned}
\]

Use distance \(+\infty\) to an empty finite set and value zero for a
maximum over an empty point set. The inequality holds because each point
in any matching either pays a real-point cost or its diagonal cost.
Nearest-neighbor queries can evaluate H efficiently.

For an increasing path, \(H(\delta)\leq C\leq A+T\). If
\(H(\delta)=A+T\), saturation is proved without solving the shifted
matching problem. With an error-controlled lower estimate \(H_{low}\)
and upper estimate \(A_{up}\), the increasing deficit is bounded by
\(\kappa_+\leq A_{up}+T-H_{low}\). This is a witness-based audit using
native diagrams and the chosen shift. It is independent of centered
matching computation, but not independent of the chosen anchor or
diagram geometry.

For a decreasing path the universal lower bound is \(C\geq A-T\). A
valid explicit shifted matching with maximum cost \(\leq A-T\) proves
equality. More generally a matching upper bound \(U\) gives
\(\kappa_-=T+C-A\leq T+U-A_{low}\). Hausdorff alone generally cannot
establish this upper bound.

Hausdorff is not an exact replacement for bottleneck matching: \(D\)
containing two copies of \((0,4)\) and \(E\) containing one has
\(H(0)=0\) but \(f(0)=2\). Multiplicity and Hall-type matching
constraints matter.

\subsection{5. Ranking, tolerance and partial
orders}\label{ranking-tolerance-and-partial-orders}

For allele a let its exactly saturated path be
\(B_a(\lambda)=A_a+\lambda k_a\), \(0\leq \lambda\leq 1\), where
\(k_a=r_a|\delta_a|\). In a numerically certified version, enclose this
line by the endpoint-derived tube. Rank crossings of exact lines occur
only where

\[
(A_a-A_b)+\lambda(k_a-k_b)=0.
\]

For a top-set rule that includes scores within \(\tau\) of the maximum,
also include the roots where this expression equals \(\pm \tau\). Retain
only roots within \([0,1]\), and handle parallel or identical lines
separately. The pairwise roots partition the path into cells with a
constant tolerance-defined top set. The boundaries require the inclusive
convention used by the frozen analysis.

An exact native tie can break for every \(\lambda>0\) if \(\tau=0\) and
slopes differ. With positive tolerance \(\tau\), that pair separates
only after \(\lambda|k_a-k_b|>\tau\). Thus top-credit changes can be
large even when distance changes are arbitrarily small; the
discontinuity belongs to the rank/tie functional, not to persistence
distance. At fixed positive \(\tau\) the discontinuities move to the
tolerance boundary.

For arbitrary certified score intervals \([L_a,U_a]\), a exceeds b by
more than \(\tau\) whenever \(L_a>U_b+\tau\). The set
\(\{a: U_a\geq \max_{b\ne a}L_b-\tau\}\) conservatively contains every
possible tolerance-defined top allele. The set
\(\{a: L_a\geq \max_{b\ne a}U_b-\tau\}\) is guaranteed to belong to the
top set for every admissible score vector. These possible and guaranteed
sets are exact for independent Cartesian score intervals; with dependent
errors they remain conservative. They define a robust partial ranking
rather than forcing unsupported total ranks. For a whole path, use the
endpoint envelopes at \(u=\lambda|\delta_a|\), or take the appropriate
infima and suprema over the \(\lambda\) interval to certify uniform
membership or pairwise ordering.

The standard anchor-blind margin guarantee is
\(A_i-A_j>|\delta_i|+|\delta_j|+\tau\), with numerical errors added when
estimated quantities are used. A slope-aware path certificate can be
tighter: minimize \((A_i-A_j)+\lambda(k_i-k_j)\) over the desired
\(\lambda\) interval, which requires only its endpoints, then subtract
the two tube errors and compare to \(\tau\). These are elementary
robust-ranking consequences, not new generic ranking theorems.

Do not choose \(\lambda\) using selected-allele labels and treat the
resulting enrichment as predeclared. Replication analyses should freeze
the tolerance, error budget, endpoint-saturation statistic, ranking
audit and any predictive baseline before map outcomes are inspected.

\subsection{6. Rank-filtration
perturbations}\label{rank-filtration-perturbations}

Assign edge ij the fraction of eligible edges with strictly greater
\(L\) value, grouping ties so equal weights enter together. The
resulting rank flag filtration is unchanged under every strictly
increasing transformation of \(L\) (Giusti et al., 2015). It removes
magnitude information more broadly than global centring. Near ties,
small entrywise changes can alter ranks, so the weighted-filtration
stability bound does not transfer automatically. If edge perturbations
are bounded by \(\delta\), only comparisons whose original gap is at
most \(2\delta\) can change. For a fixed set of \(m\) eligible edges,
the rank-filtration perturbation is at most
\(\max_e \#\{f\ne e:|w_f-w_e|\leq2\delta\}/m\). At exact ties no
universal continuity bound tending to zero with \(\delta\) exists:
splitting an all-tied four-bin graph into a zero-birth perimeter and two
arbitrarily later diagonals creates a rank interval of length 2/3.

An uncentred weighted filtration, a centred weighted filtration and a
rank filtration answer different questions. Their comparison is more
informative than labelling topology uniformly intensity-sensitive.
Degree-one interval lengths also omit genomic feature identity:
relabelling bins leaves the abstract diagram unchanged.
Coordinate-sensitive comparisons remain necessary when loop location is
the target.

\subsection{7. Verification and prior-art
scope}\label{verification-and-prior-art-scope}

\begin{itemize}
\tightlist
\item
  Sheehy, Kisielius and Cavanna, \emph{Computing the Shift-Invariant
  Bottleneck Distance for Persistence Diagrams}, CCCG 2018, already
  define optimization over diagonal shifts and study its computation.
  Their parameterized matching problem is directly relevant. The shift
  quotient and the idea of following shifted costs must not be presented
  as new: \url{https://donsheehy.net/research/cavanna18computing.pdf}
\item
  Cao, Monod, Vlontzos, Schmidtke and Kainz, \emph{Topological
  information retrieval with dilation-invariant bottleneck comparative
  measures}, published 4 July 2023, develop dilation-invariant
  comparisons and explicitly relate them to shift-invariant bottleneck
  work. This is close normalization/comparison prior art, although its
  transformation and retrieval task differ:
  \url{https://academic.oup.com/imaiai/article/12/3/1964/7217093}
\item
  May, Krishnamoorthy and Gambill, \emph{A Normalized Bottleneck
  Distance on Persistence Diagrams and Homology Preservation under
  Dimension Reduction}, submitted 11 June 2023, revised 29 August 2024
  and published in La Matematica, develops scale normalization and
  stability in a dimension-reduction setting. It further excludes broad
  novelty claims about normalized persistence comparison:
  \url{https://arxiv.org/abs/2306.06727} ;
  \url{https://doi.org/10.1007/s44007-024-00130-0}
\end{itemize}

The endpoint rigidity lemma follows from elementary Lipschitz geometry.
The finite-cell slope statement follows from finite bottleneck matching.
The interval ranking statements are standard margin arguments. No
priority claim is made for these individual results. The review covered
these closest identified sources, not every work on parametric matching
or robust ranking.

\subsection{8. Exact finite checks and their
limits}\label{exact-finite-checks-and-their-limits}

\nolinkurl{scripts/astra_verify_anchor_theory.py} independently
enumerates all partial matchings on 250 small rational diagram pairs. It
passed 498 directional local-slope cases (1,494 exact evaluations) and
4,500 endpoint-cone/deficit-tube checks, together with saturating,
V-turn, diagonal-cap, opposite-slope and multiplicity controls. Output:
\nolinkurl{outputs/astra_anchor_theory_verification.json}.

All verification arithmetic uses Python Fraction. These finite checks
support implementation and catch sign/degeneracy errors; the written
proofs carry the universal statements. They are not a proof-assistant
formalization and do not certify the full genomic diagram computations.

\end{document}
